Read--verify is the method the orchestrator runs (Figure~\ref{fig:dataflow}). It needs no training and has three
parts.

\paragraph{M1: typed abstention.}
When no reader serves a link, the method abstains with the wall type as its cause, and the model never sees the
page. When a link is served but no verified citation supports an answer after one re-ask, it abstains with cause
\texttt{unsupported}. An answer therefore takes one of three forms: a cited answer, an abstention with a wall
cause, or an abstention for lack of support.

\paragraph{M2: channel-located retrieval.}
The records of every served post (Section~\ref{sec:harness}) are ranked for the question with BM25 over
character bigrams for Chinese and lower-cased words otherwise, per channel, and fused across channels by
reciprocal rank \citep{samuel2025mmmorrf}. The twelve best records reach the model, each labelled with its
channel and locator.

\paragraph{M3: verified-quote citations with one re-ask.}
The model must return its answer with record identifiers and quotes copied from those records. Code checks every
quote against the cited record after a normalization that ignores spacing, punctuation and case. If no quote
verifies, the model is asked once more and told why; if none verifies again, the answer becomes an abstention
with cause \texttt{unsupported}. The rule is enforced in code rather than in the prompt: a prompt can ask for exact
quotes, but only code can refuse an inexact one.

Every channel that the access profile allows is extracted for every served post; the method spends no model call
on choosing what to extract.
